% ============================================================
% ch22.tex —— 第 22 章 欧拉的筛子：素数藏在分析里
% ============================================================
\chapter{欧拉的筛子：素数藏在分析里}

1737 年，欧拉拧动了数学史上最大胆的一次并网：把"素数"（离散、一颗一颗、不可预测）接上"分析"（连续、流动、可微可积）的电网。他只用了一页纸。这一章我们把那页纸走完——终点是素数定理反应堆的点火现场，也是"解析数论"这门学科的出生证明。

\section{热身：几何级数，一条万能齿轮}

\begin{kaiwei}
\[
1 + \frac12 + \frac14 + \frac18 + \frac1{16} + \cdots = ?
\]
\end{kaiwei}

第 13 章的解方程套路：设 $S$ 为总和，则 $\tfrac12 S$ 是"从第二项起"的总和。相减，中间项全部抵消：
\[
S - \frac S2 = 1 \quad\Longrightarrow\quad S = 2 .
\]
一般化（对任何 $|x| < 1$，同样的减法合法——相减时中间项错位相消）：
\[
\boxed{\;1 + x + x^2 + x^3 + \cdots = \frac{1}{1-x}\qquad (|x|<1).\;}
\]
这叫\textbf{几何级数}。$x = \tfrac13$：左边 $1 + \tfrac13 + \tfrac19+\cdots = \tfrac32$ \checkmark；$x = \tfrac12$：$= 2$ \checkmark。（若 $|x| \ge 1$——比如 $x = 1$，即调和级数——右边分母归零，机器报错，这正是"发散"的信号弹。）

请看清这个等式的\textbf{形状}：左边是\textbf{无穷连加}，右边是一个\textbf{分式}。$\tfrac1{1-x}$ 是一条翻译通道——把"加法语言"译成"除法语言"。欧拉的天才在于：他发现素数正是两条语言之间的\textbf{口岸}。

\section{欧拉的大胆一步：让素数各管一摊}

回忆算术基本定理（第 3 章）：每个正整数拆成素数幂的乘积，且\textbf{拆法唯一}——$360 = 2^3\times3^2\times5$，且全世界只有这一种配方。欧拉盯着"唯一"二字问出一个疯狂的问题：\textbf{能不能让"素数"通过乘法，把所有整数每个不重不漏地造一遍？}

给每个素数装一条几何级数"齿轮"（$x = \tfrac1p$）：
\[
\frac{1}{1-\frac12} = 1 + \frac12 + \frac14 + \frac18 + \cdots \qquad (\text{齿轮}1)
\]
\[
\frac{1}{1-\frac13} = 1 + \frac13 + \frac19 + \cdots \qquad (\text{齿轮}2)
\]
\[
\frac{1}{1-\frac15} = 1 + \frac15 + \frac1{25} + \cdots \qquad (\text{齿轮}3)
\]
现在把\textbf{所有齿轮乘起来}，用分配律展开。展开的规则：每一项都是"从每个齿轮里各挑一项、连乘"的结果。比如从齿轮 $1$ 挑 $\tfrac14$、齿轮 $2$ 挑 $\tfrac13$、齿轮 $3$ 挑 $1$（其余齿轮全挑 $1$），得到 $\tfrac14\times\tfrac13 = \tfrac1{12}$。关键在于清点：\textbf{每种挑法对应一个形如 $\tfrac{1}{2^a3^b5^c\cdots}$ 的项}，而 $2^a3^b5^c\cdots$ 恰好遍历每一个正整数（这就是"每个整数都有素数配方"），且\textbf{每个整数只被遍历一次}（这就是"配方唯一"）。先用有限个齿轮验一把（只用 $2, 3, 5$）：
\[
\left(1+\tfrac12+\tfrac14+\cdots\right)\left(1+\tfrac13+\tfrac19+\cdots\right)\left(1+\tfrac15+\cdots\right)
= 2\times\frac32\times\frac54 = 3.75,
\]
展开收集的是"仅由 $2,3,5$ 组装"的数的倒数：
\[
1 + \tfrac12 + \tfrac13 + \tfrac14 + \tfrac15 + \tfrac16 + \tfrac18 + \tfrac19 + \tfrac1{10} + \tfrac1{12} + \tfrac1{15} + \cdots
\]
（注意 $7, 11, 14$ 等缺席——它们的配方用了名单外的素数。）数值抽查：$1+0.5+0.333+0.25+0.2+0.167+0.125 = 2.575$，继续加 $\tfrac19, \tfrac1{10},\dots$ 蹭向 $3.75$——从下方逼近，方向对。把\textbf{全部}素数的齿轮乘开，展开式恰好收集\textbf{每一个正整数}的倒数，一次不多、一次不少：

\begin{dingli}[欧拉乘积公式，1737 年]
\[
\underbrace{\frac{1}{1-\frac12}\cdot\frac{1}{1-\frac13}\cdot\frac{1}{1-\frac15}\cdot\frac{1}{1-\frac17}\cdots}_{\text{乘法的世界（只认识素数）}}
\;=\;
\underbrace{1 + \frac12 + \frac13 + \frac14 + \frac15 + \cdots}_{\text{加法的世界（认识所有整数）}}
\]
简记：$\displaystyle\prod_{p}\frac{1}{1-\frac1p} = \sum_{n=1}^\infty \frac1n$（$\prod$ 是连乘记号，$p$ 跑遍全部素数）。
\end{dingli}

\textbf{左认素数，右认全体整数，中间站着唯一分解定理}——两千年前欧几里得的"原子论"，在这里现出了分析学的原形。

\section{第一次收割：素数无穷的第二证明}

立即兑现。右边是调和级数——第 21 章刚证明它\textbf{发散}（$=\infty$）。假设素数只有有限个，左边就是有限个有限数（每条齿轮 $\tfrac{1}{1-1/p}$ 都是老实的数）的乘积——有限个有限数相乘，乘积有限。有限 $=$ 无穷？矛盾。\textbf{素数有无穷多个——第二种证明，完成。}

与欧几里得的证明（第 4 章）并排欣赏：欧几里得用"乘起来加一"在\textbf{代数}里逼出新素数，构造性强（能真造出名单外的因数）；欧拉用"加法世界爆炸"在\textbf{分析}里逼出无穷，定量性强（顺便告诉你素数"质量"有多大）。同一个定理，两片大陆，两种风味——数学里好定理的标志之一，就是经得起多种证明的反复曝光。

\section{第二次收割：素数不稀少（本章主菜）}

真正的大菜。追问第 21 章末挂出的问题：
\[
\frac12 + \frac13 + \frac15 + \frac17 + \frac1{11} + \frac1{13} + \cdots \;=\; ?
\]
素数的倒数之和——收敛还是发散？这是素数密度的"终极化验"（收敛 $=$ 素数可忽略；发散 $=$ 素数撑得起一片天）。欧拉的判断：\textbf{发散}。给一个骨架完整、每步可查的论证：

\textbf{第一步（部分乘积压住调和数）。}只用不超过 $x$ 的素数造"有限齿轮组"，展开收集的是"配方只含这些素数"的数——它包含\textbf{所有 $n \le x$}（$n\le x$ 的配方当然不会用到超 $x$ 的素数）。于是
\[
\prod_{p \le x}\frac{1}{1-\frac1p} \;\ge\; \sum_{n\le x}\frac1n = H_x .
\]
（$x = 7$ 数值验算：左边 $2\times1.5\times1.25\times\tfrac76 = 4.375$，右边 $H_7 = 2.593$，压住 \checkmark。）

\textbf{第二步（取对数，连乘变连加）。}两边取 $\ln$（第 19 章法则：$\ln$ 把乘法拆成加法）：
\[
\sum_{p\le x}\ln\frac{1}{1-\frac1p} \;\ge\; \ln H_x .
\]
由第 21 章，$H_x \approx \ln x + \gamma \to \infty$，故右边 $\to\infty$。

\textbf{第三步（每一项与 $\tfrac1p$ 同量级）。}检查单个齿轮的对数：用几何级数展开 $\ln\dfrac1{1-u}$（对 $\tfrac1{1-u}$ 的级数逐项积分，或直接认可）：
\[
\ln\frac{1}{1-u} = u + \frac{u^2}{2} + \frac{u^3}{3} + \cdots \;\le\; u + u^2 + u^3 + \cdots = \frac{u}{1-u} \;\le\; 2u \quad (u \le \tfrac12).
\]
取 $u = \tfrac1p$（素数最小是 $2$，$u \le \tfrac12$ 恰好达标）：
\[
\ln\frac{1}{1-\frac1p} \;\le\; \frac{2}{p},
\qquad\text{从而}\qquad
\frac1p \;\ge\; \frac12\,\ln\frac{1}{1-\frac1p}.
\]

\textbf{第四步（合流）。}把第三步代入第二步：
\[
\sum_{p\le x}\frac1p \;\ge\; \frac12\sum_{p\le x}\ln\frac{1}{1-\frac1p} \;\ge\; \frac12 \ln H_x \;\longrightarrow\; \infty .
\]
\[
\boxed{\;\frac12 + \frac13 + \frac15 + \frac17 + \frac1{11} + \cdots = \infty\;}
\]
\textbf{素数的倒数之和发散}——素数没有稀到可以忽略不计。而且发散的\textbf{速度}也可用同款机器量出（精修版，梅尔滕斯定理 1874）：
\[
\sum_{p\le x}\frac1p \;\approx\; \ln\ln x + 0.2615\cdots
\]
（常数 $0.2615$ 叫梅尔滕斯常数。）$\ln\ln x$——对数套对数，慢得离谱（$x = 10^{18}$ 时 $\ln\ln x \approx 3.78$：人类数过的所有素数的倒数加起来还不到 $4$）——但确凿发散。这条"慢而无穷"的曲线，正是素数密度 $\tfrac1{\ln x}$ 的镜像：密度按 $\tfrac1{\ln x}$ 稀释，倒数的累积按 $\ln\ln$ 爬升，\textbf{两者是同一枚硬币的两面}（$\int\tfrac{dt}{t\ln t} = \ln\ln t$——第 21 章烧脑时刻你亲手算过的积分）。素数定理的全部图景，已经在这一步的阴影里了。

\begin{cidian}
土名"齿轮"$\to$ \textbf{欧拉乘积}；把指数 $1$ 换成变量 $s$ 的升级版 $\displaystyle\sum_n \frac{1}{n^s} = \prod_p \frac{1}{1-p^{-s}}$，左边叫\textbf{$\zeta$ 函数}（读"截塔"），$\zeta(s)$——第 24 章的主角，黎曼猜想的舞台。
\end{cidian}

\begin{lishi}
莱昂哈德·欧拉（1707--1783），史上最高产的数学家（全集八十余卷，生前失明后靠心算与口述续写了近半），"读起来像随风写字的人"。欧拉乘积、$\mathrm{e}$ 的命名、$\ln$ 记号、$\sum$ 记号、$\mathrm{e}^{\mathrm{i}\pi}+1=0$、巴塞尔问题……本书已经欠他太多人情，而第 23、24 章还要继续欠。他 1737 年写下这条乘积公式时未必料到：一百六十年后，黎曼会把它搬进复平面，撬动整个素数王国。
\end{lishi}

\begin{dongshou}
\begin{itemize}
  \yisi 数值验证齿轮组：只用 $2, 3, 5, 7$ 四个齿轮，算乘积 $2\times\tfrac32\times\tfrac54\times\tfrac76 = 4.375$；再算 $H_{10} = 2.929$，验证"部分乘积 $\ge H_x$"。然后加上齿轮 $11$（$\tfrac{1}{1-1/11} = 1.1$），乘积跳到 $4.8125$，观察"每加一个齿轮，乘积乘上略大于 $1$ 的数"——齿轮的增量正是素数稀疏的代数化身。
  \yier 用几何级数公式求和：$1 + \tfrac13 + \tfrac19 + \cdots$（$\tfrac32$）；$1 + \tfrac1{10} + \tfrac1{100}+\cdots$（$\tfrac{10}{9}$）；并解释 $x = 1$ 时机器为何报错（分母 $1-x = 0$，对应调和级数发散——通道两头是通的）。
  \yier 手算素数倒数和到 $p = 29$：$\tfrac12+\tfrac13+\tfrac15+\tfrac17+\tfrac1{11}+\tfrac1{13}+\tfrac1{17}+\tfrac1{19}+\tfrac1{23}+\tfrac1{29}$（$\approx 1.354$）。对照精修公式 $\ln\ln 29 + 0.2615 \approx \ln(3.367) + 0.2615 \approx 1.214 + 0.262 = 1.476$——小数据偏差正常（误差项 $O(1/\ln x)$），体会"渐近公式越晚越准"的含义。
  \yisi 不计算，只查本章表格判断：$\tfrac12+\tfrac13+\cdots+\tfrac1{997}$（$168$ 个素数）大约是几？用 $\ln\ln 997 + 0.2615 \approx \ln(6.9) + 0.26 \approx 2.16$。真值 $2.107$——三分钟前你还不认识这个数，现在能估它。
\end{itemize}
\end{dongshou}

\begin{shaonao}
\textbf{把欧拉的机器再榨一遍：素数密度的粗刻度。}从齿轮的对数里不但能证明发散，还能读出密度。

(1) 对 $\ln\tfrac1{1-u} = u + \tfrac{u^2}{2} + \tfrac{u^3}{3}+\cdots$，当 $u$ 小（大素数）时首项主导：$\ln\tfrac1{1-1/p} \approx \tfrac1p$。于是第二步的不等式近似变成等式方向：
\[
\sum_{p\le x}\frac1p \approx \sum_{p\le x}\ln\frac1{1-1/p} \approx \ln\prod_{p\le x}\frac1{1-1/p} \approx \ln H_x \approx \ln\ln x .
\]

(2) 反过来问：若 $x$ 附近素数密度是 $f(x)$，则素数倒数和约为 $\int^x \dfrac{f(t)\,dt}{t}$（每个素数 $p$ 贡献 $\tfrac1p$，密度加权）。令它等于 $\ln\ln x$：
\[
\int^x\frac{f(t)}{t}\,dt = \ln\ln x = \int^x\frac{dt}{t\ln t} \quad\Longrightarrow\quad f(t) = \frac{1}{\ln t} .
\]
——\textbf{素数密度 $\tfrac1{\ln t}$ 从欧拉的齿轮里反解了出来}。第 20 章的"筛子直觉"在此兑现为推导。你已经用初等工具摸到了素数定理的门口（门口与门内还差一步——第 23 章的"误差控制"，那一步需要 $\zeta$ 的复分析）。
\end{shaonao}

\begin{chengshi}
诚实的三处地毯。其一，"分配律展开\textbf{无穷}个\textbf{无穷}和"在欧拉时代是自由操作，十九世纪分析严格化后需要担保——严格做法是先证 $\sum \tfrac1{n^s}$ 在 $s>1$ 时绝对收敛（对 $s=1$ 即调和级数，发散），对 $s>1$ 做严格的欧拉乘积，再令 $s \downarrow 1$ 取极限，本章为叙事流畅直接在 $s = 1$ 处展开，方向与结论均正确，细节是标准的"取极限"练习。其二，第三步的 $\ln\tfrac1{1-u}$ 级数来自逐项积分，引用不证。其三，"$\sum_{p\le x}\tfrac1p \approx \ln\ln x + 0.2615$"是梅尔滕斯定理，其误差项与黎曼猜想直接挂钩（若 RH 成立，误差可压到 $O(\ln x)$ 级）——地毯下面还是地毯，这正是数学的层叠结构。
\end{chengshi}

反应堆点火成功：齿轮咬出了密度 $\tfrac1{\ln x}$。下一章，读出仪表盘上的正式读数——\textbf{素数定理}，以及它身后那台若隐若现的复分析机器。
